\section{An infinite cyclic image}\label{sec:image}

We now compute the image of the surface constructed in
Proposition~\ref{prop:orbifold-surface}.  The relations supplied by a
coupling hyperplane show that this image is cyclic.  Proving that the
cyclic group is infinite requires a separate argument: the translation
labels of the preceding section exist in charts, and need not define a
homomorphism on the fundamental group of the whole orbifold.

\begin{theorem}\label{thm:orbifold-cyclic}
For a sufficiently deep choice of the arithmetic level in the
construction of Proposition~\ref{prop:orbifold-surface}, the map
$f:Y\longrightarrow\cS$ satisfies
\[
 \im\bigl(f_*:\pi_1(Y)\longrightarrow\pi_1(\cS)\bigr)\cong\Z.
\]
Here $Y$ is an iterated point blowup of the simple abelian surface
$\Jac(C)$, and its image lies in the ordinary locus of $\cS$.
\end{theorem}

Throughout the proof, fundamental groups are based at points in the
ordinary locus.  Changes of basepoint are made along specified common
paths; in particular, the comparisons below concern homomorphisms with
one common conjugation, rather than separate conjugations of individual
generators.

\subsection{Surface loops and the cyclic upper bound}

Recall the distinguished fiber $F_*=C\times C$ in $V$, the order-four
lift $s$, and the rational section used to construct $Y$ in
Section~\ref{geo:surface}.  Put
\[
 \mathcal K=[V/\langle s\rangle].
\]
There is a natural homomorphism
\begin{equation}\label{app:comparison-map}
 \pi_1(\mathcal K)\longrightarrow\pi_1(\cS).
\end{equation}
Indeed, the substack $[F_*/\langle s\rangle]$ has complex codimension
two in the smooth stack $\mathcal K$.  Removing it does not change the
orbifold fundamental group, and its complement is the open substack on
which the construction of $\cS$ leaves $\mathcal K$ unchanged.

Choose a general point $q_0\in C$.  The map
\[
 C\longrightarrow\Sym^2 C\longrightarrow\Jac(C),
 \qquad p\longmapsto p+q_0,
\]
is, up to translation, an Abel map.  Its map on fundamental groups
surjects onto $\pi_1(\Jac(C))=H_1(C,\Z)$.  Since point blowups preserve
fundamental groups, loops from $C\times\{q_0\}$ supply generators of
$\pi_1(Y)$.  They may be represented by a finite bouquet avoiding the
branch points, $p=q_0$, the zeros and poles of the rational section,
and the finitely many points where the rational map or its resolution
must be avoided.  Removing finitely many points from a smooth curve is
surjective on fundamental groups, so these restrictions do not lose any
generators.

\begin{lemma}\label{app:surface-comparison}
With common choices of basing, the images in $\pi_1(\cS)$ of the
preceding generators of $\pi_1(Y)$ equal the images, under
\eqref{app:comparison-map}, of the corresponding loops in
$C\times\{q_0\}\subset V$.  The loops can also be represented in the
sign-covered arrangement complement, arbitrarily close to $F_*$, with
both normal coordinates nonzero.
\end{lemma}

\begin{proof}
Lift the bouquet by the rational section to the exceptional divisor of
$\operatorname{Bl}_{F_*}V$.  Over this bouquet the section has the form
\[
 [\tau(p),\sqrt{-1}\,\tau(q_0)],
\]
with both entries nonzero.  It avoids the fixed curves involved in the
last resolution step, because the bouquet avoids the diagonal of
$C\times C$.

A complex line bundle on a finite graph is topologically trivial.
Thus the normal line to this exceptional divisor admits one continuous
nonvanishing section over the entire lifted bouquet.  In the product normal coordinates, this push can be written
explicitly as
\[
 \varepsilon\bigl(\tau(p),\sqrt{-1}\,\tau(q_0)\bigr).
\]
One sufficiently small $\varepsilon>0$ works over the compact bouquet
and gives a closed based bouquet off the exceptional divisor.  On blowing down to $V$, the pushed bouquet is
homotopic to the original one in $F_*$.  On passing to the resolution
neighborhood in $\cS$, it is homotopic to the bouquet obtained from
$Y$.  These homotopies use the same push of the common basepoint.

The two components of the displayed section are nonzero, and the
chosen directions avoid the marked slopes.  Taking the push small
therefore puts its projection off the arrangement.  The distinguished
point is an isolated component of the fixed locus of every nonidentity
power of $\sigma$ in a sufficiently small neighborhood.  Hence the
pushed bouquet can also avoid the inverse images of all those fixed
loci.  This proves all the assertions, including the common basing.
\end{proof}

We next compute an upper bound before taking the quotient by $s$.
Write $H=(\Z/2)^2$ for the sign deck group.  The map
\[
 V\longrightarrow[V/H]
\]
is an orbifold covering, so it induces an injection of $\pi_1(V)$ into
$\pi_1([V/H])$.  This remains true over branch points: the target is the
quotient stack, which retains their isotropy groups.

The pencil through $L_z$ gives the orbifold line
$[C/(\Z/2)]$, with six order-two points.  Number the four red points
first.  Its fundamental group has the presentation
\begin{equation}\label{app:pencil-group}
 Q=\langle h_1,\ldots,h_6\mid
       h_i^2=1,\ h_1h_2h_3h_4h_5h_6=1\rangle.
\end{equation}
The kernel of the parity character sending every $h_i$ to $1\in\Z/2$
is $\pi_1(C)$.  With the $w$ sign point fixed over a general direction,
the map of this orbifold pencil to $[V/H]$ identifies the homomorphism
from this even subgroup with the homomorphism from
$\pi_1(C\times\{q_0\})$ to $\pi_1(V)$, followed by the injection just
described.

\begin{lemma}\label{app:cyclic-upper}
The image of $\pi_1(C\times\{q_0\})$ in $\pi_1(V)$ is cyclic.
Consequently the image of $\pi_1(Y)$ in $\pi_1(\cS)$ is cyclic.
\end{lemma}

\begin{proof}
We show that the four red meridians in \eqref{app:pencil-group} have
the same image in $\pi_1([V/H])$.  First move the $w$ coordinate a little
off its exceptional direction line, keeping it fixed, generic, and
small.  Push the red pencil meridians off the $z$ exceptional divisor.
They can be represented with
\[
 t=z_1\ne0\quad\hbox{fixed and small},\qquad
 \lambda=z_2/z_1
\]
and with the red meridians based by paths in a small disk containing
exactly the four red values of $\lambda$.

Slide this calculation in the $t$ variable to the coupling hyperplane
$t=c$.  The disk of red directions can be chosen small enough that
this entire motion, with the chosen small $w$, remains inside the
ball: the red slopes have modulus less than $0.01$, whereas
$c/\sqrt a<0.9$.  The motion and the finitely many paths used here lie
in a fixed compact part of the full finite model.  We include this
compact set among the regions that must agree with the actual
arrangement when the arithmetic level is chosen.

Before blowing up the pair intersections, the relevant local divisor
is
\[
 \{t=c\}\ \cup\ \bigcup_{i=1}^4\{\lambda=\lambda_i\}.
\]
Let $k$ be the meridian in the $t$ variable, with one common choice of
basing.  The product structure makes $k$ commute with each red
meridian.  More precisely, the same $t$ circle can be transported along
each of the based paths in the direction disk; this gives one element
$k$, not four unrelated conjugates.

Blowing up $\{t=c,\lambda=\lambda_i\}$ introduces an exceptional
divisor whose meridian is $kh_i$: in the two transverse coordinates,
the exceptional meridian circles both coordinates once.  The filling
rules of Section~\ref{geo:fill} impose no branching along this
exceptional divisor, since two branch components met there.  Thus its
meridian is killed.  Therefore
\[
 kh_i=1\qquad(1\le i\le4).
\]
The strict transforms retain order two, so $k^2=h_i^2=1$.  All four
red images are consequently one involution, denoted $r$.

The product relation in \eqref{app:pencil-group} now gives
$h_5h_6=1$, and the two blue images are one involution $b$.  The image
of $Q$ is thus a quotient of
\[
 \langle r,b\mid r^2=b^2=1\rangle.
\]
Every even word in $r,b$ is a power of $rb$.  Hence the image of the
even subgroup $\pi_1(C)$ is cyclic.  Injectivity of
$\pi_1(V)\to\pi_1([V/H])$ gives the assertion in $\pi_1(V)$.
Lemma~\ref{app:surface-comparison} and the surjectivity from the chosen
$C$ loops onto $\pi_1(Y)$ give the final assertion.
\end{proof}

This argument allows the cyclic image to be finite.  We will exclude
that possibility by applying Lemma~\ref{thm:path-lemma} to paths in a
dense open subset of $\cS$.

\subsection{Paths in the ordinary open set}

Let $O\subset\cS$ be the sign-covered arrangement complement,
quotiented by $s$, after also removing the full inverse image of the
fixed loci in $M$ of all nonidentity powers of $\sigma$.  The
modifications defining $\cS$ are isomorphisms over this set.  In
particular $O$ is a dense connected smooth variety, and all quotient
actions used over it are free.

There is a useful space for specifying lifts.  Over the complement in
$\BB$ of the lifted arrangement and the lifted omitted fixed loci,
pull back the sign data from $V$, and call the resulting space
$O^\#$.  Then
\[
 O^\#\longrightarrow O
\]
is an ordinary covering.  Changes of lift are generated by the ball
deck group $\Gamma$ and by $\sigma$, the latter acting on the sign data
through $s$.  These are precisely the transformations used to form
$O$; we do not divide out the two sign deck transformations.  The group generated by $\Gamma$ and $\sigma$ acts freely on the
retained ball region.

Fix the diameter bound $u$ supplied by
Proposition~\ref{prop:model-hyperbolicity}.  Define a graph $G$ whose
vertices are the points of $O$.  An edge is a parametrized continuous
path in $O$ whose lifted projection to $\BB$ has diameter strictly
less than $u$.  All graphs use the same parametrization and reversal
conventions.  The diameter condition is independent of the lift,
because changes of lift act by ball isometries.  No bound is imposed
on path length or winding.  In particular, arbitrarily high powers of
a loop can be individual edges when their projections stay in one
small region.

\subsection{Choosing and comparing the charts}

We make the order of the parameters explicit.  First fix $u$, then a
common hyperbolicity constant for the model graphs, and then the
integer $N$ required by Lemma~\ref{thm:path-lemma}.  Set
\begin{equation}\label{app:radius}
 R=(2N+100)(u+1).
\end{equation}
A lift of a graph path with at most $N$ edges stays within distance
$Nu$ of its initial ball point, regardless of the lengths of the
individual parametrized paths.

The special points of $\BB$ are the points of $\Gamma o$, the lifts
of the distinguished point of $M$.  Choose $J>10R$ so large that, in
the central model, both of the following hold at every point of
distance at least $J-3R$ from $o$:
\begin{enumerate}
\item On sign loops projected into the $R$ ball about that point, the
kernel of the sum translation equals the intersection of the two
individual translation kernels.
\item For $j=1,2,3$, the displacement by $\sigma^j$ exceeds $10R$.
\end{enumerate}
The first requirement follows from
Proposition~\ref{prop:kernel-compatibility}.  For the second, none of
$\sigma,\sigma^2,\sigma^3$ has a boundary fixed point.  The ball
distance formula then implies that its displacement tends uniformly
to infinity as the viewpoint tends to the boundary.

Let $B(R)$ be the buffer radius in
Proposition~\ref{prop:kernel-compatibility}, enlarged so that
$B(R)>R$, and choose
\begin{equation}\label{app:chart-radius}
 K>10\bigl(J+R+B(R)\bigr).
\end{equation}
All these constants depend only on the finite models.  Now take the
initial arithmetic level deep enough that
Proposition~\ref{prop:arrangement-local} holds on balls of radius
$10K$, that such balls inject into the torsion-free ball quotient,
that the full finite arrangement agrees with the lifted arrangement
in these neighborhoods of every special point, and that distinct
special points have distance greater than $100K$.  Enlarge this last
choice, if necessary, to include the fixed compact paths used in
Lemma~\ref{app:cyclic-upper}.  Subsequent finite covers preserve all
these requirements.

For a vertex $v\in G$, choose a lift with ball coordinate $x$ and its
actual sign data.  Its chart has one of two forms.
\begin{description}
\item[Central chart.] If $\dist(x,\Gamma o)<J$, use the full finite
model centered at the unique nearby special point.  Read in the
sum-translation cover and then quotient by the lifted order-four
symmetry, as in Lemma~\ref{mod:central}.
\item[Individual chart.] Otherwise take all lifted hyperplanes
meeting $B_{\BB}(x,K)$, identify them, with their prescribed normal
vectors, with the corresponding subset of $\cD$, and use the model
cover given by the two individual translation kernels.  There is no
symmetry quotient in this chart.
\end{description}
These chart graphs are the graphs of
Proposition~\ref{prop:model-hyperbolicity}, with the appropriate
additional analytic omissions.  One can transport the whole locally
finite union of lifted omitted fixed loci to model coordinates; any
extra lifted arrangement hyperplanes outside the comparison region
may be omitted as well.  Such omissions are globally defined on the
model ball.  In central coordinates they are invariant under the
symmetry.  Proposition~\ref{prop:model-hyperbolicity} shows that these
omissions do not enlarge the uniform hyperbolicity constant.

For central charts, fix the identifications of sign data once on the
large comparison ball about $o$.  Sign characters on that ball
complement are determined by their meridians.  Choose the sheet
identifications to agree with the local construction near $F_*$;
they then intertwine the actual lift $s$ with the model lift
throughout this connected region.  Transport the choices to all
special points by $\Gamma$.  These choices are also compatible with
$\sigma$, by the symmetry of the construction.

Within the comparison ball about a special point, the only quotient
transformations that can identify two ball points are the four
symmetry powers centered there.  Indeed, if $g\in\langle\Gamma,
\sigma\rangle$ identifies two such points, it carries the special
point to another one at distance at most twice the comparison radius.
Separation forces it to fix that special point.  Its stabilizer is
exactly the corresponding conjugate of $\langle\sigma\rangle$,
since $\Gamma$ is torsion free.

To read a path of at most $N$ graph edges from $v$, lift it from the
chosen point of $O^\#$.  Its ball projection remains in the comparison
region.  Identify the resulting sign path with the model sign path,
lift further from a chosen initial point in the translation cover,
and, in a central chart, take the class modulo the lifted symmetry.
This is a path in a model graph $D_v$, with initial state $d_v$.
Write $[p]_v$ for its terminal state.

Central models are centered only at points of $\Gamma o$.  Other
full pencil fibers, including those over
$\Gamma_0o\setminus\Gamma o$, and the other quotient fixed loci
remain unmodified in $[V/\langle s\rangle]$.  An ordinary-open loop around such a locus may have a ball or
sign lift with different endpoints; an individual chart retains this
difference.  Equality of actual endpoints is not required to imply
equality of read states.  For orbifold relations, the relevant
meridian power lifts to a closed loop in a smooth chart in $V$;
its closure will be checked separately below.

\begin{lemma}\label{app:chart-axioms}
The readings just constructed satisfy the three chart hypotheses of
Lemma~\ref{thm:path-lemma}.
\end{lemma}

\begin{proof}
The uniform hyperbolicity assertion is
Proposition~\ref{prop:model-hyperbolicity}.  We verify the two path
compatibility assertions.

\emph{Continuation from a state.}
Path lifting is consistent with prefixes and reversals.  Equal read
states project to the same actual point of $O$.  In an individual
chart, equality means equality in the model covering, and hence also
in the ball and sign data.  Every outgoing model edge at a state
reached after fewer than $N$ edges gives an actual edge of $G$: its
projection stays within $u$ of its starting ball point and therefore
inside the comparison region.  It avoids precisely the arrangement
and extra loci excluded in $O$.  Lifting and projecting parametrized
paths are mutually inverse once the starting lift is fixed.  This is
the required bijection on outgoing edges, and it depends only on the
read state.

In a central chart the symmetry acts freely on the open complement.
An outgoing edge of its quotient graph consequently has a unique
representative from any chosen representative of the starting
vertex.  Representatives at equal states differ by a symmetry power,
which respects the identification with the actual open set.  Thus the
same outgoing-edge bijection holds, including its dependence only on
the state.  This proves the first chart hypothesis.

\emph{Comparison after a prefix.}
Let $p$ go from $v$ to a vertex $v'$, and let $a,b$ start at $v'$,
with
\[
 |p|,|a|,|b|\le N/2.
\]
We must prove
\begin{equation}\label{app:prefix-compatibility}
 [pa]_v=[pb]_v\quad\Longleftrightarrow\quad[a]_{v'}=[b]_{v'}.
\end{equation}
For this calculation, use for $v'$ the lift reached by lifting $p$.
The equality tests do not depend on this change from its initially
chosen lift.  For $\Gamma$ this follows from transport of the normal
data.  For $\sigma$ the two families are interchanged and the finite
arrangement has the same symmetry.  If relabeling changes a chosen
normal vector by a scalar, it changes the corresponding translated
normal by the same scalar, because $\sigma$ normalizes the
arithmetic group.  Thus the intrinsic prescriptions of
Proposition~\ref{prop:kernel-compatibility} give the same individual
kernels, with possible reversals of translation orientations.
Individual sign-sheet choices likewise do not affect a kernel test
on a closed sign loop.  For central charts the fixed, equivariant
identifications above give the same conclusion for the sum cover.
Finally, the choice of initial point in a translation cover is
irrelevant because the covers are regular, and the deck translations
commute with the lifted central symmetry.

All projected paths in the comparison are contained in the $R$ ball
about the endpoint of the lifted $p$.  If both charts are individual,
equality first requires coincident endpoints in the actual ball and
sign data.  If this requirement fails, both tests fail.  If it holds,
compare the two sign paths by the loop $a\overline b$.  The further
lifts end together exactly when both translations of this loop
vanish.  The common prefix does not change that condition.  Both
charts contain the required buffered ball about this new viewpoint,
by \eqref{app:chart-radius}; therefore
Proposition~\ref{prop:kernel-compatibility} identifies the two kernel
tests.  This proves \eqref{app:prefix-compatibility} in this case.

If both charts are central, their special points coincide.  Otherwise
these two special points would be at distance at most $2J+Nu$,
contrary to their separation.  The two readings then use restrictions
of the same covering with the same lifted symmetry, so their endpoint
equality tests agree.

It remains to compare a central and an individual chart.  The
viewpoint at the end of $p$ is at distance at least $J-R$ from the
special point of the central chart.  This follows directly from the
criterion for the individual root and from the bound $Nu<R$ on the
distance between the two roots.  Every endpoint being compared is
therefore at distance at least $J-2R$ from that special point.  Two
such endpoints in the comparison ball cannot differ by a nontrivial
symmetry power: their mutual distance is at most $2R$, whereas the
relevant displacement exceeds $10R$.  Thus the central test, too,
requires equality of the ball and sign endpoints.  On the remaining
sign loop, the sum kernel equals the intersection of the individual
kernels by the choice of $J$.  The buffered comparison of individual
kernels now applies as before.  This proves
\eqref{app:prefix-compatibility} in the mixed case, and completes the
verification of the chart hypotheses.
\end{proof}

\subsection{From orbifold relations to path relations}

Let $\mathcal P(G)$ be the path groupoid of $G$ modulo backtracking
and every edge loop of at most three edges that closes in its starting
chart.  Lemma~\ref{thm:path-lemma} and
Lemma~\ref{app:chart-axioms} detect nontrivial elements of this
groupoid.  To use that conclusion for $\cS$, we need the following
direction of comparison: every loop that is trivial in
$\pi_1(\cS)$ must already be trivial in $\mathcal P(G)$.

We recall explicitly the topological fact about orbifolds used in
this comparison.

\begin{lemma}\label{app:orbifold-meridians}
Let $\mathcal T$ be a connected smooth effective complex orbifold, and
let $U$ be a dense ordinary open subset whose complement is complex
analytic.  Then $\pi_1(U)\to\pi_1(\mathcal T)$ is surjective.  Its
kernel is normally generated by the loops obtained from boundaries
of small transverse disks in smooth orbifold charts at general
points of the omitted divisors.  Equivalently, if the generic
stabilizer along such a divisor has order $m$, one kills the $m$th
power of a small meridian in the ordinary quotient.  Omitted strata
of complex codimension at least two add no relations.
\end{lemma}

\begin{proof}
One may apply ordinary general position to a smooth frame
presentation of $\mathcal T$.  Effectiveness and finite-group
linearization imply that the action of each chart group on tangent
frames is free.  Consequently the frame space is a manifold, with a
locally free action of the connected group $\operatorname{GL}_n(\C)$,
and its quotient stack is $\mathcal T$.  The homotopy fibration from
this presentation, as in \cite[Example~5.6]{Noohi2014}, computes $\pi_1(\mathcal T)$ as the fundamental
group of the frame space modulo the image of
$\pi_1(\operatorname{GL}_n(\C))$; the same description applies over
$U$.

In the frame space, paths can avoid real-codimension-two loci, and a
nullhomotopy disk can meet such loci transversely in finitely many
general divisor points.  Removing small disks around those
intersection points expresses its boundary as a product of
conjugates of transverse meridians.  Real codimension at least four
can be avoided by the entire disk.  Passing through the preceding
quotient of fundamental groups proves the assertion.  At a general
point of a divisor, the effective stabilizer is cyclic and acts
faithfully on the transverse line.  A circle in that line projects
to the $m$th power of a quotient meridian, giving the last description.
\end{proof}

\begin{lemma}\label{app:true-relations}
A loop of edges of $G$ that is nullhomotopic in $\cS$ is trivial in
$\mathcal P(G)$.
\end{lemma}

\begin{proof}
First consider homotopies within $O$.  Every ordinary path admits a
finite subdivision into graph edges.  An already specified edge may
itself be subdivided without changing its class in $\mathcal P(G)$:
if $e=e_1e_2$ is a subdivision, then $e_1,e_2$ are still edges, and
$e_1e_2\overline e$ is a three-edge loop that closes on reading,
since all three pieces lift the same parametrized path.  Repeating
this observation allows arbitrary finite subdivision even when an
original edge winds arbitrarily often.

A homotopy in $O$ can now be subdivided into sufficiently small path
triangles.  To choose the neighborhoods uniformly over its compact
image, take lifted ball patches avoiding the arrangement and the
omitted fixed loci, each disjoint from its distinct quotient
translates.  For every root in a smaller patch, the oversized chart
agrees with this same arrangement-free ball patch.  The sign and
translation covers trivialize there, and a central symmetry quotient
makes no additional identifications inside the patch.  A finite
cover of the homotopy image by these smaller neighborhoods supplies
the required subdivision.  Each triangle consequently closes in its
starting chart and is one of the imposed relations.  This is the
usual edge-path proof that ordinary homotopy relations are respected.

It remains, by Lemma~\ref{app:orbifold-meridians}, to check the
transverse-disk relations.  Away from the distinguished fiber,
$\cS$ is $[V/\langle s\rangle]$.  Even at points with inertia, the
required power of a quotient meridian lifts to a closed small loop
bounding a disk in the smooth local uniformizing space in $V$.
Its projection to $M$ bounds a disk in an arbitrarily small
simply connected ball patch, so its further lift to $\BB$ is closed
and lies in the corresponding small region.  The finite model there
has exactly the same participating hyperplanes and the same blowups, sign branching, and fillings as
$V$.  This local identification extends across the fillings by the
explicit local rules, or equivalently by uniqueness of the
normalizations extending the sign covers.  Lemma~\ref{mod:extension}
therefore says that the translation covers extend over this disk.
Its boundary closes in an individual chart, and also closes in a
central chart by the sum rule.  This argument includes unbranched
exceptional divisors, order-two branch divisors, and the powers
required by any further quotient inertia.  Additional analytic loci
removed in defining $O$ cause no new obstruction: the extension test
is made after those additional omissions are restored.

Over the distinguished point, the relevant chart is central.  By
Lemma~\ref{mod:central}, its sum cover descends to the coarse quotient
and pulls back to a covering of the resolution neighborhood.  Thus
it extends over every transverse disk used for the ordinary smooth
space there.  Those meridians close in the central graph as well.
A sufficiently small representative is a single edge: its ball lift
stays arbitrarily close to the distinguished point, even when its
endpoint differs from its start by a central symmetry power.
Away from this fiber the same diameter assertion follows from the
small local chart in $V$, including for a power of a meridian.

Every required transverse-disk boundary is consequently a
single-edge loop that closes in its chart, and is killed in
$\mathcal P(G)$.  Normal generation, together with the homotopy
argument in $O$, proves the lemma.
\end{proof}

Combining the last two lemmas with Lemma~\ref{thm:path-lemma} gives
\begin{equation}\label{app:detection}
 \text{a single edge loop with distinct read endpoints represents a
 nonidentity element of }\pi_1(\cS).
\end{equation}
This conclusion does not require any of the translation labels to
extend to all of $\pi_1(\cS)$.

\subsection{Infinite order and completion of the proof}

Choose a based loop $\alpha$ in $C\times\{q_0\}$ on which the first
translation label is nonzero.  Such a loop comes from lifting a
product of a red and a blue meridian in the six-point pencil: its
parity is even and its dihedral label is $rb$, a nontrivial
translation.  The second label is zero on this factor.  The loop
can be represented in the good locus used in
Lemma~\ref{app:surface-comparison}, since deleting its finite
exceptional set is surjective on $\pi_1(C)$ and the translation label
extends over the filled curve.

Push $\alpha$ to a loop $e$ in $O$ as in that lemma, so close to the
distinguished fiber that its ball projection has diameter less than
$u$.  Its basepoint has a central chart.  The sum label on the pushed
loop is the same nonzero integer as on $\alpha$, by extension of the
sum covering over the filled model.  For every nonzero integer $n$,
the parametrized loop $e^n$ has the same projected image as $e$, and
hence is itself one edge of $G$.

The read endpoints of $e^n$ are distinct in the central quotient
graph.  Before taking the symmetry quotient they differ by the
nonzero translation $n\,\ell(\alpha)$.  The lift of $e$ is closed in
the actual ball and sign data, because the push is closed upstairs in $V$ and its projection lies
inside the prescribed small simply connected ball about $o$.  If the two translated endpoints became
equal after quotienting by a lifted symmetry power, that power would
fix their common ball point.  This point lies in the ordinary open
set where the symmetry acts freely, so the power would have to be
the identity.  A nonzero deck translation cannot fix a point of a
covering.  Thus the endpoints remain distinct.

By \eqref{app:detection}, every nonzero power of the class of $e$ is
nontrivial in $\pi_1(\cS)$.  Lemma~\ref{app:surface-comparison} puts
this infinite-order element in the image of $\pi_1(Y)$.  That image
is cyclic by Lemma~\ref{app:cyclic-upper}; it is therefore infinite
cyclic.  This proves Theorem~\ref{thm:orbifold-cyclic}.
